% AY2012 制御工学 解答例
% 回答の流れ・言葉遣いは 03_制御工学/参考/「制御工学Ⅱ（完成版）」に寄せる
% （align の ∴ チェーンで淡々と計算、簡潔な言い回し）。デザイン（囲み枠等）は真似しない。
% 参考はあくまで書き方の手本。解答・数値・問題は本年度過去問から自力で導いた。
\documentclass[11pt,a4paper]{article}
\usepackage{../../../00_common/preamble}

% --- 作図用（preamble に未定義のため本ファイルで補う）---
\usetikzlibrary{positioning,arrows.meta,calc}
\usepackage{pgfplots}
\pgfplotsset{compat=newest}

\title{制御工学 AY2012 解答例（専門応用科目I）}
\author{}
\date{}

\begin{document}
\maketitle

% ==================================================================
\section*{第1問}

二段階制御システム（内側ループ：$G_1(s)$ を $K_0$ で負帰還、
外側ループ：$G_2(s)$ を単位負帰還）について答える。

\subsection*{(1) 伝達関数}
まず内側ループ（$K_0$ 負帰還）を閉じる。$G_2$ の出力を $W$ とすると
\begin{align*}
  &\qquad Y = G_1\,(W - K_0 Y)\\
  &\therefore\quad (1+G_1K_0)\,Y = G_1 W
   \quad\therefore\quad M(s)\equiv\frac{Y}{W}=\frac{G_1}{1+G_1K_0}.
\end{align*}
次に外側ループ（単位負帰還）を閉じる。$W=G_2(U-Y)$ なので
\begin{align*}
  &\qquad Y = M\,G_2\,(U-Y)\\
  &\therefore\quad (1+M G_2)\,Y = M G_2\,U\\
  &\therefore\quad \frac{Y(s)}{U(s)}=\frac{M G_2}{1+M G_2}
   =\frac{G_1G_2}{1+G_1K_0+G_1G_2}.
\end{align*}
\[
  \therefore\quad \ans{\;\dfrac{Y(s)}{U(s)}=\dfrac{G_1(s)\,G_2(s)}{1+G_1(s)K_0+G_1(s)G_2(s)}\;}
\]

\subsection*{(2) $G_1(s)$ のゲイン曲線}
$G_1(s)=\dfrac{1}{(s+1)(10s+1)}$、$K_0=1$ とする。ゲインは
\begin{align*}
  20\log_{10}|G_1(j\omega)|
  &= -20\log_{10}\sqrt{1+\omega^2}\;-\;20\log_{10}\sqrt{1+(10\omega)^2}.
\end{align*}
直流ゲインは $|G_1(0)|=1$（$0\,\mathrm{dB}$）、折れ点は $10s+1$ より $\omega=0.1$、$s+1$ より $\omega=1\,[\mathrm{rad/s}]$。
傾きは $\omega<0.1$ で $0$、$0.1<\omega<1$ で $-20$、$\omega>1$ で $-40\,[\mathrm{dB/dec}]$ となる。

\begin{center}
\begin{tikzpicture}
\begin{axis}[
  width=0.86\linewidth, height=5.4cm,
  xmin=1e-2, xmax=1e2, ymin=-120, ymax=20, xmode=log,
  xlabel={$\omega\,[\mathrm{rad/s}]$}, ylabel={ゲイン [dB]},
  ytick={-120,-100,-80,-60,-40,-20,0,20},
  xmajorgrids=true, ymajorgrids=true, xminorgrids=true,
  grid style={line width=0.3pt, draw=gray!60},
  extra y ticks={0}, extra y tick style={grid=major, grid style={line width=1pt, draw=black}},
  legend pos=south west, clip=true, enlargelimits=false,
]
\addplot[thick] coordinates {(0.01,0) (0.1,0) (1,-20) (100,-100)};
\addlegendentry{漸近線}
\addplot[thick,dashed] coordinates
  {(0.01,0) (0.1,-3) (0.316,-13) (1,-23) (10,-60.1) (100,-100)};
\addlegendentry{真値}
\node[anchor=south] at (axis cs:0.1,0)  {\scriptsize $\omega=0.1$};
\node[anchor=north east] at (axis cs:1,-20) {\scriptsize $\omega=1$};
\end{axis}
\end{tikzpicture}
\end{center}

\[
  \therefore\quad \ans{\;0\,\mathrm{dB}\ (\omega\le0.1)\ \xrightarrow{-20\,\mathrm{dB/dec}}\ \omega=1\ \xrightarrow{-40\,\mathrm{dB/dec}}\;}
\]

\subsection*{(3) $G_2=K_1$ のときの定常値}
$G_1(0)=1$、$K_0=1$、$G_2=K_1$ を (1) に代入する。系は安定なので、単位ステップ入力に対して
最終値の定理より
\begin{align*}
  &\qquad y(\infty)=\lim_{s\to0}sY(s)=\lim_{s\to0}T(s)
   =\frac{G_1(0)K_1}{1+G_1(0)K_0+G_1(0)K_1}\\
  &\therefore\quad y(\infty)=\frac{K_1}{1+1+K_1}=\frac{K_1}{K_1+2}.
\end{align*}
\[
  \therefore\quad \ans{\;y(\infty)=\dfrac{K_1}{K_1+2}\;}
\]
$K_1$ が有限である限り定常偏差 $1-y(\infty)=\dfrac{2}{K_1+2}$ が残る。

\subsection*{(4) $G_2(s)=\dfrac{s+0.1}{s}$ のときの定常値}
$G_2(s)=\dfrac{s+0.1}{s}=1+\dfrac{0.1}{s}$ は積分動作（原点に極）をもつ。単位ステップ入力に対し
最終値の定理より
\begin{align*}
  &\qquad y(\infty)=\lim_{s\to0}T(s)
   =\lim_{s\to0}\frac{G_1G_2}{1+G_1K_0+G_1G_2}\\
  &\therefore\quad s\to0\ \text{で}\ G_2\to\infty\ \text{なので、分母・分子とも}\ G_1G_2\ \text{が支配的となり}\\
  &\therefore\quad y(\infty)=\lim_{s\to0}\frac{G_1G_2}{G_1G_2}=1.
\end{align*}
なお特性方程式 $10s^3+11s^2+3s+0.1=0$ は係数がすべて正かつ $11\cdot3>10\cdot0.1$ より安定だから、
最終値の定理は適用できる。
\[
  \therefore\quad \ans{\;y(\infty)=1\;}\qquad(\text{定常偏差}\ 0)
\]

\subsection*{(5) 補償器 $G_2(s)$ の役割}
$G_2(s)=\dfrac{s+0.1}{s}$ は原点の極（積分器 $1/s$）と $s=-0.1$ の零点をもつ
PI（比例積分）型の直列補償器である。役割は次の二点である。
\begin{itemize}\setlength{\itemsep}{1pt}
  \item 積分動作によって系の型が1上がり、ステップ入力に対する定常偏差を $0$ にする
        （(3) の偏差 $\tfrac{2}{K_1+2}$ が、(4) では $0$ になった）。
  \item 零点 $s=-0.1$ が積分器の位相遅れ $(-90^\circ)$ を補い、位相余裕を確保して安定性を保つ。
\end{itemize}
\[
  \therefore\quad \ans{\;\text{積分動作で定常偏差を $0$ にし、零点で位相余裕を確保する PI 型直列補償器}\;}
\]

\bigskip
\noindent 検算: (1) は内側・外側ループ別導出と直接展開が一致（SymPy、差 $=0$）。
(2) は $\omega=100$ の真値 $-100\,\mathrm{dB}$ が漸近線と一致。(4) は特性方程式の 3 根の実部がいずれも負（安定）✓。

\newpage
% ==================================================================
\section*{第2問}

一自由度振動系（前向き $\frac{1}{M}\to\frac1s\to\frac1s$、速度帰還 $D$・位置帰還 $K$）について答える。

\subsection*{(1) 伝達関数}
加速度を $A$、速度を $V$ とすると、ブロック線図より
\begin{align*}
  &\qquad A=\frac{1}{M}\bigl[(U-KY)-DV\bigr],\quad V=\frac1s A,\quad Y=\frac1s V.\\
  &\therefore\quad V=sY,\ A=s^2Y\ \text{を代入して}\quad Ms^2Y=U-KY-DsY\\
  &\therefore\quad \bigl(Ms^2+Ds+K\bigr)Y=U.
\end{align*}
\[
  \therefore\quad \ans{\;\dfrac{Y(s)}{U(s)}=\dfrac{1}{Ms^2+Ds+K}\;}
\]

\subsection*{(2) 固有角振動数と減衰係数}
標準形に直して、$s^2+2\zeta\omega_n s+\omega_n^2$ と係数を比較する。
\begin{align*}
  &\qquad \frac{Y}{U}=\frac{1/M}{s^2+\dfrac{D}{M}s+\dfrac{K}{M}}\\
  &\therefore\quad \omega_n^2=\frac{K}{M},\qquad 2\zeta\omega_n=\frac{D}{M}.
\end{align*}
\[
  \therefore\quad \ans{\;\omega_n=\sqrt{\dfrac{K}{M}}\;},\qquad
  \ans{\;\zeta=\dfrac{D}{2\sqrt{KM}}\;}
\]

\subsection*{(3) インパルス応答（$M=1,\ D=3,\ K=2$）}
パラメータを代入し、インパルス入力のラプラス変換は $F(s)=1$ なので
\begin{align*}
  &\qquad G(s)=\frac{1}{s^2+3s+2}=\frac{1}{(s+1)(s+2)}\\
  &\therefore\quad X(s)=G(s)\,F(s)=\frac{1}{s+1}-\frac{1}{s+2}\\
  &\therefore\quad x(t)=\mathcal{L}^{-1}\!\left[X(s)\right]=\e^{-t}-\e^{-2t}.
\end{align*}
\[
  \therefore\quad \ans{\;h(t)=\e^{-t}-\e^{-2t}\quad(t\ge0)\;}
\]

\subsection*{(4) 単位ステップ応答}
ステップ入力のラプラス変換は $F(s)=\dfrac1s$ なので
\begin{align*}
  &\qquad X(s)=G(s)\,F(s)=\frac{1}{s(s+1)(s+2)}
   =\frac{\frac12}{s}+\frac{-1}{s+1}+\frac{\frac12}{s+2}\\
  &\therefore\quad x(t)=\mathcal{L}^{-1}\!\left[X(s)\right]
   =\frac12-\e^{-t}+\frac12\e^{-2t}.
\end{align*}
\[
  \therefore\quad \ans{\;y(t)=\dfrac12-\e^{-t}+\dfrac12\e^{-2t}\quad(t\ge0)\;}
\]

\bigskip
\noindent 検算: (1) を SymPy で照合。(3) の極は $s=-1,-2$（過減衰 $\zeta=\tfrac{3}{2\sqrt2}>1$）。
(4) の定常値 $y(\infty)=\tfrac12=G(0)=\tfrac1K$、初期値 $y(0)=0,\ \dot y(0)=0$ を確認✓。

\end{document}
